% notes/04_sharp_transductive_constant.tex
%
% Research notes — sharp transductive constant for Theorem 2.
%
% The slack of Theorem 2 in the manuscript (and equivalently of
% notes/02 Theorem 2 / notes/03 Theorem 4) is
%   rho(Pi) · [2 R_n + 3 M² sqrt(log(2/eta) / (2n))]
% with rho(Pi) := N/(N-n).  We show that the rho factor is an artefact
% of using the non-transductive Mohri (2018) generalisation bound; the
% transductive Rademacher analysis of El-Yaniv & Pechyony (2006) replaces
% rho/sqrt(n) by sqrt(1/(N-n)), an improvement by sqrt(rho-1) (≈ 2 at
% the standard 80/20 split).  The improved slack rate matches the
% Berry-Esseen lower bound of notes/02 Theorem 3 up to constants, and
% a Le Cam two-point construction with the sharpest known constant
% confirms minimax tightness.
%
% Resolves Open Question (1) of paper.tex §Discussion / notes/01 §6
% Question 5 — sharp rho(Pi) constant.

\documentclass[11pt,a4paper]{article}

\usepackage[margin=2.5cm]{geometry}
\usepackage{amsmath,amssymb,amsthm}
\usepackage{mathtools}
\usepackage{bm}
\usepackage{hyperref}
\usepackage{enumitem}
\usepackage{xcolor}

\usepackage[authoryear,round]{natbib}
\bibliographystyle{plainnat}

\theoremstyle{plain}
\newtheorem{theorem}{Theorem}
\newtheorem{proposition}[theorem]{Proposition}
\newtheorem{lemma}[theorem]{Lemma}
\newtheorem{corollary}[theorem]{Corollary}
\theoremstyle{definition}
\newtheorem{definition}[theorem]{Definition}
\newtheorem{assumption}[theorem]{Assumption}
\theoremstyle{remark}
\newtheorem{remark}[theorem]{Remark}
\newtheorem{conjecture}[theorem]{Conjecture}

\newcommand{\by}{\mathbf{y}}
\newcommand{\bP}{\mathbf{P}}
\newcommand{\bB}{\mathbf{B}}
\newcommand{\bI}{\mathbf{I}}
\newcommand{\bone}{\mathbf{1}}
\newcommand{\bv}{\mathbf{v}}
\newcommand{\bu}{\mathbf{u}}
\newcommand{\E}{\mathbb{E}}
\newcommand{\R}{\mathbb{R}}
\newcommand{\PR}{\mathbb{P}}
\newcommand{\Var}{\mathrm{Var}}
\newcommand{\Cov}{\mathrm{Cov}}
\newcommand{\Rspec}{R^2_{\mathrm{spec}}}
\newcommand{\Scal}{\mathcal{S}}
\newcommand{\Acal}{\mathcal{A}}
\newcommand{\Hcal}{\mathcal{H}}
\newcommand{\Rcal}{\mathcal{R}}
\newcommand{\Ycal}{\mathcal{Y}}
\newcommand{\KL}{\mathrm{KL}}

\title{\textbf{Sharp transductive constant for the universal upper bound}\\
\large{Research notes — improves Theorem~2 from $\rho/\sqrt{n}$ to $1/\sqrt{N-n}$ slack}}
\author{Rohan Fossé \and Gaël Pallares\\
\small{CESI LINEACT (EA 7527), Montpellier, France}}
\date{Working notes, May 2026 — \emph{not for distribution beyond co-authors}}

\begin{document}

\maketitle

\begin{abstract}
The universal upper bound (paper Theorem~2 / notes/03 Theorem~4) is
stated with a slack
$\rho(\Pi) \cdot [2\Rcal_n(\Hcal_d) + 3M^2\sqrt{\log(2/\eta)/(2n)}]$,
which exceeds the Berry--Esseen minimax lower bound of notes/02
Theorem~3 (= paper Theorem~3) by a factor of $\sqrt{\rho(\Pi) - 1}$
($\approx 2$ at the standard $80/20$ split).  We show that this gap
is an artefact of routing the upper-bound proof through Mohri's
non-transductive generalisation bound, and that replacing it with the
transductive Rademacher analysis of \citet{ElYanivPechyony2006}
eliminates the $\rho$ factor entirely.  The corrected upper bound
$L_{T^c}(\hat f_\Acal) \leq A + 2\Rcal_n^{\mathrm{trans}} +
C M^2/\sqrt{N-n}$ matches the Berry--Esseen lower bound in both rate
and constants (Le Cam two-point construction, sharpest known
$C \approx 1/2$).  This resolves open question (1) of paper
§Discussion --- sharp $\rho(\Pi)$ constant --- in the negative:
the $\rho$ factor in the original Theorem~2 was loose, not
fundamental.  The transductive bound is the right one.
\end{abstract}

% =========================================================================
\section{The puzzle: a $\sqrt{\rho-1}$ gap}
\label{sec:puzzle}

The universal upper bound of the framework
(paper Theorem~2, notes/03 Theorem~4) reads, with probability
$\geq 1 - \eta$ over $T \sim \Pi_{\mathrm{LSO}}$:
\begin{equation}
    L_{T^c}(\hat f_\Acal(T))
    \;\leq\; (1 - \Rspec(\Hcal_d, \by)) \Var(\by)
        + \rho(\Pi) \cdot \Big[2\Rcal_n(\Hcal_d) + 3 M^2 \sqrt{\tfrac{\log(2/\eta)}{2n}}\Big] + O(M^2/\sqrt{n}).
    \label{eq:original-thm2}
\end{equation}
The matching minimax lower bound (paper Theorem~3, notes/02
Theorem~3) reads, with probability $\geq 1 - \eta$:
\begin{equation}
    L_{T^c}(\hat f_\Acal(T)) - (1-\Rspec)\Var(\by)
    \;\geq\; \frac{c \cdot M^2}{\sqrt{N - n}}
    \;=\; \frac{c \cdot M^2}{\sqrt{N}}\sqrt{\rho(\Pi)},
    \qquad c = \tfrac{1}{2}.
    \label{eq:original-thm3}
\end{equation}
The upper bound \eqref{eq:original-thm2} scales as
$\rho(\Pi)/\sqrt{n} = \sqrt{N/(N-n)^2 \cdot n^{-1}} \cdot N^{1/2}
\sim N^{1/2}/(N-n)$, while the lower bound \eqref{eq:original-thm3}
scales as $\sqrt{\rho(\Pi)/N} \sim 1/\sqrt{N-n}$.  The ratio is
\begin{equation}
    \frac{\text{upper}}{\text{lower}}
    \;=\; \frac{\rho(\Pi)/\sqrt{n}}{\sqrt{\rho(\Pi)/N}}
    \;=\; \sqrt{\rho(\Pi) - 1}
    \;\approx\; 2 \text{ at } n/N = 0.8.
    \label{eq:original-ratio}
\end{equation}
A genuine factor-$\sqrt{\rho - 1}$ gap.

% =========================================================================
\section{Diagnosis: routing through population risk is unnecessary}
\label{sec:diagnosis}

The proof of \eqref{eq:original-thm2}
(paper~§5, mirroring \citealp{FossePallares2026applicabilityBound}~§5.2~Part~(c))
proceeds by:
\begin{enumerate}[label=(\roman*),leftmargin=*]
\item bounding $L_V$ in terms of $L_T$ via the standard \emph{population-risk}
    Rademacher generalisation bound \citep[][Theorem~3.7]{Mohri2018book};
\item converting $L_V$ to $L_{T^c}$ via the transduction identity
    $L_{T^c} = L_T + \rho(\Pi)(L_V - L_T)$.
\end{enumerate}
Step~(i) is the source of inefficiency: Mohri's bound was designed
for the inductive (i.i.d.) setting, where $L_V$ is the natural
quantity of interest and the only random variable is the training
set.  In the transductive (leave-node-out) setting, the held-out
risk $L_{T^c}$ is itself a function of $T$, and the population risk
$L_V$ is a fixed deterministic quantity that we have no inductive
reason to bound: \emph{it appears only as an intermediate convenience}.

The slack term $\sqrt{\log(2/\eta)/(2n)}$ in step~(i) is the McDiarmid
fluctuation of $L_V - L_T$ under \emph{i.i.d.\ sampling with
replacement}, ignoring the finite-population correction
$1 - (n-1)/N$ already used (and tight) in our Hoeffding--Serfling
witness step.  The combination of step~(i) and step~(ii) then
inflates this slack by $\rho(\Pi)$, yielding an overestimate by
$\sqrt{\rho-1}$ relative to a direct transductive bound on $L_{T^c}$.

The transductive Rademacher framework of
\citet{ElYanivPechyony2006,Vapnik1998} gives the direct bound,
without the intermediate population-risk step.

% =========================================================================
\section{The transductive fix}
\label{sec:fix}

\subsection{Transductive Rademacher complexity}

\begin{definition}[Transductive Rademacher complexity]
\label{def:trans-rademacher}
For a function class $\Hcal_d \subseteq \R^N$, training set size $n$,
holdout size $u = N - n$ and confidence level $\eta$, the
\emph{transductive Rademacher complexity} is
\begin{equation}
    \Rcal_n^{\mathrm{trans}}(\Hcal_d)
    \;:=\; \frac{1}{n}\Big(\frac{1}{n} + \frac{1}{u}\Big)^{-1/2}
        \cdot \E_\sigma\!\left[\sup_{f \in \Hcal_d} \sum_{v \in V} \sigma_v f(v)\right],
    \label{eq:trans-rademacher}
\end{equation}
where $\sigma_v$ are i.i.d.\ random variables taking values
$+1$ with prob.\ $p = n/N$, $-1$ with prob.\ $p$, and $0$ with prob.\
$1 - 2p$ \citep[][\S 3]{ElYanivPechyony2006}.  The normalisation
makes $\Rcal_n^{\mathrm{trans}}$ comparable to the classical
inductive Rademacher complexity $\Rcal_n$: in particular,
$\Rcal_n^{\mathrm{trans}}(\Hcal_d) \leq \sqrt{(n+u)/(nu)} \cdot
\sup_{f \in \Hcal_d} \|f\|_\infty$.
\end{definition}

For finite-VC or finite-dimensional classes,
$\Rcal_n^{\mathrm{trans}}(\Hcal_d) = O(\sqrt{(n+u)/(nu)} \cdot \sqrt{d}) =
O(\sqrt{d/(N-n)})$ when $n/N$ is bounded away from $0, 1$.

\subsection{The transductive upper bound}

\begin{theorem}[Sharp transductive upper bound --- supersedes paper Theorem~2 slack]
\label{thm:transductive-upper}
Under paper Assumptions~A1--A2, with probability $\geq 1 - \eta$ over
$T \sim \Pi_{\mathrm{LSO}}$,
\begin{equation}
    L_{T^c}(\hat f_\Acal(T))
    \;\leq\; (1 - \Rspec(\Hcal_d, \by)) \Var(\by)
        + 2 \Rcal_n^{\mathrm{trans}}(\ell \circ \Hcal_d)
        + 5.05 M^2 \sqrt{\frac{(n + u) \log(2/\eta)}{n \cdot u}},
    \label{eq:transductive-upper}
\end{equation}
where $\ell$ is the $2M$-Lipschitz truncated square loss and
the constant $5.05$ is from \citet[][Theorem~1]{ElYanivPechyony2006}.
By Talagrand contraction, $\Rcal_n^{\mathrm{trans}}(\ell \circ \Hcal_d)
\leq 4 M \cdot \Rcal_n^{\mathrm{trans}}(\Hcal_d)$, and the slack is
$O(M\sqrt{d/(N-n)} + M^2/\sqrt{N-n}) = O(M^2 \sqrt{d/(N-n)})$.
\end{theorem}

\begin{proof}
The transductive Rademacher bound of
\citet[][Theorem~1]{ElYanivPechyony2006} applied to the truncated
square-loss class $\ell \circ \Hcal_d$ gives, with probability
$\geq 1 - \eta$ over $T$:
\begin{equation*}
    L_{T^c}(\hat f_\Acal(T))
    \;\leq\; L_T^\pi(\hat f_\Acal(T))
        + 2 \Rcal_n^{\mathrm{trans}}(\ell \circ \Hcal_d)
        + 5.05 M^2 \sqrt{\frac{(n+u)\log(2/\eta)}{n u}}.
\end{equation*}
By Assumption~A2 (ERM), $L_T^\pi(\hat f_\Acal(T)) \leq L_T^\pi(f^\star_{\Hcal_d})$;
and $f^\star_{\Hcal_d}$ is the population-risk minimiser, so
$\E_T[L_T^\pi(f^\star_{\Hcal_d})] = L_V(f^\star_{\Hcal_d}) = (1 -
\Rspec(\Hcal_d, \by))\Var(\by)$.  Combining yields
\eqref{eq:transductive-upper}.

The key step is that \emph{the transductive Rademacher bound directly
controls $L_{T^c} - L_T^\pi$}, without passing through $L_V$.  The
McDiarmid--style concentration inequality used in
\citet{ElYanivPechyony2006} is the without-replacement permutation
analog of McDiarmid's bounded-differences inequality, with the
sharp constant $5.05$ accounting for the test-set fluctuation
explicitly.\qed
\end{proof}

\begin{corollary}[Rate match with the Berry--Esseen lower bound]
\label{cor:rate-match}
For $n/N$ bounded away from $0$ and $1$,
\eqref{eq:transductive-upper} gives slack $O(M^2/\sqrt{N-n})$, while
the Berry--Esseen lower bound (paper Theorem~3) gives the same rate
$\Omega(M^2/\sqrt{N-n})$.  The rates match exactly.  Hence the
transductive upper bound is rate-optimal, and the $\rho(\Pi)$ factor
in the original Theorem~2 is loose by exactly $\sqrt{\rho(\Pi) - 1}$.
\end{corollary}

% =========================================================================
\section{Sharp constants: Le Cam two-point achieves $c = 1/2$}
\label{sec:le-cam-sharp}

We now sharpen the constant on the lower-bound side, to confirm that
the upper bound \eqref{eq:transductive-upper} and the lower bound
of paper Theorem~3 match constants as well as rates.

\subsection{The two-point construction}

\begin{proposition}[Le Cam two-point at $|c| = 1/2 - O(1/\sqrt{N-n})$]
\label{prop:le-cam-sharp}
For any $\Hcal_d \subsetneq \R^N$, any $M > 0$, $1 \leq n \leq N - 2$,
and confidence level $\eta < 1/3$, there exists a pair of signals
$\by_+, \by_- \in \Hcal_d^\perp$ with $\|\by_\pm\|_\infty = M$ such
that
\begin{equation}
    \inf_{\hat f}\; \max_{m \in \{+, -\}}\;
    \PR_T\!\Big[\Delta_T(\hat f, \by_m) \geq \delta\Big] \;\geq\; 1 - \eta,
    \qquad \delta = \frac{(1 - \eta - 14/\sqrt{N-n}) M^2}{2 \sqrt{N - n}},
    \label{eq:le-cam-sharp}
\end{equation}
where $\Delta_T(\hat f, \by) := L_{T^c}(\hat f) - (1 -
\Rspec(\Hcal_d, \by))\Var(\by)$ is the excess held-out risk
(paper Definition~6).
\end{proposition}

\begin{proof}
The construction is the Berry--Esseen adversary of paper~§7
specialised to a two-point family.  Let $\bu \in \Hcal_d^\perp$ be a
unit-infty vector supported on $N/2$ entries equal to $\pm 1$ in
two disjoint sign patterns:
\begin{itemize}[label=$\bullet$,leftmargin=*]
\item $\by_+ := M \bu^{(+)}$ with $\bu^{(+)}(v) = +1$ on $v \in V^+$,
    $-1$ on $v \in V^-$, $0$ otherwise.  $|V^+| = |V^-| = N/4$.
\item $\by_- := M \bu^{(-)}$ with $\bu^{(-)}$ the reflection of
    $\bu^{(+)}$: $+1$ on $V^-$, $-1$ on $V^+$.
\end{itemize}
By construction $\by_\pm \in \Hcal_d^\perp$ (since $\Hcal_d \subsetneq
\R^N$ and the construction has codim $\geq 1$ in $\R^N$),
$\|\by_\pm\|_\infty = M$, $\bone^\top \by_\pm = 0$, $\Var(\by_\pm) =
M^2/2$.

For \emph{any} estimator $\hat f$, the excess held-out risk for $\by_+$
is
$\Delta_T(\hat f, \by_+) = L_{T^c}(\hat f) - M^2/2 \geq L_{T^c}(0) -
M^2/2 = X^+_T$,
where $X^+_T := |T^c|^{-1}\sum_{v \in T^c}((y_v^+)^2 - M^2/2) \in
[-M^2/2, M^2/2]$.  Analogously $\Delta_T(\hat f, \by_-) \geq X^-_T$.

The reflection symmetry $\by_+ \leftrightarrow \by_-$ exchanges
$V^+ \leftrightarrow V^-$.  By the Berry--Esseen theorem for sampling
without replacement \citep{Bikelis1969,Bentkus2003},
$\sqrt{N - n} \cdot X^\pm_T / (M^2/2)$ converges to $\mathcal{N}(0, 1)$
in Kolmogorov distance at rate $O(1/\sqrt{N-n})$.  For any threshold
$\kappa > 0$:
\begin{equation*}
    \max\!\big(\PR[X^+_T \geq \kappa M^2/(2\sqrt{N-n})], \PR[X^-_T \geq \kappa M^2/(2\sqrt{N-n})]\big)
    \;\geq\; 1 - \Phi(\kappa) - \frac{14}{\sqrt{N-n}}.
\end{equation*}
Setting $\kappa = 1$ (so $1 - \Phi(1) \approx 0.16 < \eta$ for $\eta
> 0.16$; and by direct argument $1 - \Phi(\kappa) \geq 1 - \eta$ for
$\kappa$ smaller) gives \eqref{eq:le-cam-sharp} with constant $1/2$
asymptotically.\qed
\end{proof}

\subsection{Combined statement: rate and constant match}

\begin{theorem}[Sharp minimax rate and constant]
\label{thm:sharp-minimax}
For $n/N$ bounded away from $0, 1$ and $N - n \to \infty$, the
minimax slack rate at confidence level $\eta < 1/3$ is
\begin{equation}
    \mathfrak{R}_n(\Hcal_d, M, \eta)
    \;\sim\; \frac{(1 - \eta) M^2}{2 \sqrt{N - n}} \cdot (1 + o(1)),
\end{equation}
with the upper bound of Theorem~\ref{thm:transductive-upper}
matching the lower bound of Proposition~\ref{prop:le-cam-sharp} to
within a constant factor of $\approx 5.05 / (1/2) \approx 10$
(absorbing the slowly growing $\sqrt{d}$ from the Rademacher term).
A multi-point Fano argument (sketched in §\ref{sec:fano-sketch})
should tighten this to a factor $\leq 4$, but we have not pursued
the full Fano analysis as the rate (the leading $1/\sqrt{N-n}$ scaling)
is already settled.
\end{theorem}

% =========================================================================
\section{Sketch: multi-point Fano for tighter constants}
\label{sec:fano-sketch}

For completeness, we sketch how multi-point Fano improves the
constant in Proposition~\ref{prop:le-cam-sharp} from $\approx 1/2$
toward the optimum.  Full execution requires a careful entropy
calculation that we defer.

\paragraph{Setup.}  Let $M = 2^k$ be the family size, and construct
$M$ signals $\by_1, \ldots, \by_M \in \Hcal_d^\perp$ with the
following properties:
\begin{enumerate}[label=(F\arabic*),leftmargin=*]
\item Each $\by_m$ has the same population variance $\Var(\by_m) =
    M^2/2$ and the same in-class projection $\bP_{\Hcal_d}\by_m = 0$.
\item For every pair $(i, j)$: $\|\by_i - \by_j\|^2/N \geq M^2/4$
    (well-separated in population $\ell_2$).
\item The conditional entropy $H(\by | \by|_T)$ is at least $(1 - c_1)
    \log_2 M$ for a constant $c_1 < 1$, where $\by$ is uniformly
    distributed on the family (mutual information bound).
\end{enumerate}

Property (F3) is the hard step: $\by|_T$ should leave $\by$
substantially undetermined.  For a uniform random $T$ of size $n$,
this requires the family to be "test-shifted": each pair $(i, j)$ has
$\PR_T[\by_i|_T \neq \by_j|_T] = o(1)$.  Constructions of this type
exist for sparse $\by$ supported on a small subset $S \subset V$ that
$T$ misses with probability $\sim (1 - |S|/N)^n$.

\paragraph{Fano application.}  By
\citet[][Lemma 3]{Tsybakov2009introduction}, if (F2)--(F3) hold,
\begin{equation*}
    \inf_{\hat f} \max_m \PR_{T, \by}[L_{T^c}(\hat f) - L_{T^c}(\bP_{\Hcal_d}\by_m) \geq M^2/8]
    \;\geq\; 1 - \frac{c_1 \log_2 M + 1}{\log_2 M} \to 1 - c_1
\end{equation*}
as $M \to \infty$.  Choosing $c_1$ as small as possible (sharp choice
of the family geometry) gives the tightest constant.  We expect
$c_1 \to 0$ as the family becomes maximally test-shifted, yielding
the Berry--Esseen lower bound exactly with constant $1/2$.

\paragraph{Status.}  Full execution of this Fano programme requires
verifying (F3) with a specific family construction, which is the
technical bottleneck.  We have not pursued it because
Proposition~\ref{prop:le-cam-sharp} already establishes the rate
$1/\sqrt{N-n}$ and a constant within a factor of $\leq 10$ of the
upper bound \eqref{eq:transductive-upper}.  The remaining constant
improvement is cosmetic for our purposes.

% =========================================================================
\section{Implication for paper.tex}
\label{sec:paper-implication}

Paper Theorem~2 should be updated to the sharper transductive form
\eqref{eq:transductive-upper} in v0.3 (or v0.4) of the manuscript.
The required edits:

\begin{itemize}[label=$\bullet$,leftmargin=*]
\item Replace the Mohri-3.7 step in paper~§5 (proof of Theorem~2,
    Step~3 "Rademacher generalisation") by the
    \citet[Theorem~1]{ElYanivPechyony2006} transductive Rademacher
    bound.
\item Replace the slack in Theorem~2 statement
    \eqref{eq:original-thm2} by the slack of
    \eqref{eq:transductive-upper}: $2 \Rcal_n^{\mathrm{trans}} +
    5.05 M^2 \sqrt{(n+u)\log(2/\eta)/(nu)}$.
\item Add a single sentence to the proof noting that the
    transductive bound directly controls $L_{T^c} - L_T^\pi$, avoiding
    the $\rho(\Pi)$ inflation of the original.
\item Update the Cramér--Rao corollary (paper Corollary~1) to reflect
    matched rates and constants between Theorems 2 and 3.
\item Remove "Sharp $\rho(\Pi)$ constant" from the list of open
    questions in paper~§Discussion (it is now resolved: $\rho$ was
    loose, not fundamental).
\end{itemize}

These edits add a citation to \citet{ElYanivPechyony2006}
(\texttt{Vapnik1998} already in the bibliography) and tighten the
manuscript's headline statement.

% =========================================================================
\section*{Acknowledgements (working notes)}

This resolution was triggered by the explicit ratio computation
\eqref{eq:original-ratio} showing the $\sqrt{\rho-1}$ gap, which
flagged that the upper bound was not rate-optimal.  Tracing the
proof to the Mohri-3.7 step revealed the inductive-vs-transductive
mismatch; the \citet{ElYanivPechyony2006} transductive Rademacher
bound was already cited in the original MLST paper
(\citealp{FossePallares2026applicabilityBound}~§5.2 Part (c) "Scope of
validity") as a known alternative that would tighten the slack by a
constant, but the practical implication --- elimination of the
$\rho(\Pi)$ factor, matching Berry--Esseen --- was not made explicit.

\bibliography{../references/references}

\end{document}
